%% Correction audit de veracite - 2026-09-20 : citations de rapport de jury sans
%% support textuel dans le corpus (rapports 2018-2025) remplacees par des verbatims
%% verifies ou supprimees. Detail ci-dessous.
%% l.328 (ex) : seul le noyau chiffre est authentique (RJ2024 l.509 : 70 %% / 47 %%, mais
%%   la partie n'y est pas nommee "partie E (AOP)"). Les ajouts ("montages complexes
%%   (conversion T/I, decalage)", "calculs encadres avec des unites") sont INTROUVABLES
%%   -> citation reduite au verbatim RJ2024 l.509.
%% l.329 (ex) : citation "jury 2024" sur les trois conditions de l'AOP ideal :
%%   INTROUVABLE (aucun commentaire de jury sur l'AOP dans le corpus) -> remplacee par
%%   le verbatim RJ2024 l.515 sur l'exigence de justification.
\expandafter\ifx\csname CAPETINCLUDEMODE\endcsname\relax
\documentclass[11pt,a4paper]{article}
% preamble.tex — packages et commandes communes a toutes les fiches CAPET SII IE
\input{glyphtounicode}
\pdfgentounicode=1
\usepackage{cmap}
\usepackage[T1]{fontenc}
\usepackage[utf8]{inputenc}
\usepackage[french]{babel}
% babel[french] active a \begin{document} des raccourcis de ponctuation (:;!?)
% qui cassent la syntaxe to[...] de circuitikz (erreur "+ or - expected") ;
% on les desactive, ce qui ne change que l'espacement fin avant ces signes.
\AtBeginDocument{\shorthandoff{:;!?}}
\usepackage[a4paper, top=1.6cm, bottom=1.6cm, left=1.4cm, right=1.4cm]{geometry}
\setlength{\headheight}{14.5pt}
\usepackage{amsmath, amssymb, amsfonts}
\usepackage{siunitx}
\sisetup{locale=FR, output-decimal-marker={,}, inter-unit-product=\ensuremath{{}\cdot{}},
         group-minimum-digits=4}
% Unites non fournies par siunitx, declarees ici une fois pour toutes (fiches + exercices).
% Les gardes \ifdefined evitent tout conflit avec une redeclaration locale restee dans un
% fichier : ne pas retirer les gardes, ne pas redeclarer ces unites dans les fichiers.
\DeclareSIUnit{\jour}{j}
\ifdefined\var\else\DeclareSIUnit{\var}{var}\fi   % puissance reactive
\ifdefined\vad\else\DeclareSIUnit{\vad}{VAD}\fi   % puissance deformante
\ifdefined\tr\else\DeclareSIUnit{\tr}{tr}\fi      % tour (tr/min)
\usepackage{array, booktabs, multirow, tabularx}
\usepackage{xcolor}
\usepackage[most, breakable]{tcolorbox}
\usepackage{multicol}
\usepackage{enumitem}
\usepackage{mathtools}
\usepackage{lmodern}
% Refonte visuelle 2026 (direction "moderne coloré") : Lato pour le texte,
% maths sans-serif assorties via newtxsf (charge apres amsmath/amssymb).
\usepackage[default]{lato}
\usepackage{newtxsf}
\usepackage{microtype}
\usepackage{fancyhdr}
\usepackage{tikz}
\usetikzlibrary{positioning,arrows.meta}
\usepackage{circuitikz}
\usepackage{lastpage}

% ── Couleurs par domaine (palette modernisée 2026) ────────────────────
\definecolor{colElec}{RGB}{29, 78, 216}        % Électrotechnique — bleu vif
\definecolor{colElecLight}{RGB}{219, 234, 254} % fond bleu clair
\definecolor{colAuto}{RGB}{21, 128, 61}        % Automatique — vert
\definecolor{colMeca}{RGB}{194, 65, 12}        % Mécanique — orange
\definecolor{colTherm}{RGB}{185, 28, 28}       % Thermique — rouge
\definecolor{colQT}{RGB}{232, 244, 233}        % (hérité, plus utilisé par questiontype)
\definecolor{colJury}{RGB}{255, 248, 225}      % fond jury (crème)
\definecolor{colJuryFrame}{RGB}{202, 138, 4}   % bordure/bandeau jury (ambre)
\definecolor{colGrayLight}{RGB}{245, 245, 247} % fond gris neutre
\definecolor{colSeparator}{RGB}{180, 180, 180} % gris séparateur
\definecolor{colExo}{RGB}{8, 145, 178}          % Banque d'exercices — cyan (distinct des fiches)
% ── Encre et gris de la refonte ───────────────────────────────────────
\definecolor{colInk}{RGB}{23, 32, 42}          % texte courant (quasi-noir)
\definecolor{colSlate}{RGB}{83, 98, 118}       % gris ardoise (méta, en-têtes)
\definecolor{colRule}{RGB}{203, 213, 225}      % filets clairs
\definecolor{colJuryBg}{RGB}{255, 248, 225}    % fond crème jury
\definecolor{colPresence}{RGB}{83, 98, 118}    % bloc 5 (présence dans les sujets)
\AtBeginDocument{\color{colInk}}

% ── Style encadré exercice + corrigé (banque d'exercices) ─────────────
\tcbset{
  exoBox/.style={
    breakable,
    enhanced,
    colback=colExo!8!white,
    colframe=colExo,
    colbacktitle=colExo,
    coltitle=white,
    fonttitle=\bfseries\small,
    left=4pt, right=4pt, top=3pt, bottom=3pt,
    before skip=5pt, after skip=5pt,
  }
}

% ── Bascule enonce / corrige (banque d'exercices) ─────────────────────
% \CAPETCORRIGE est defini en ligne de commande (pdflatex "\def\CAPETCORRIGE{1}\input{...}")
% pour compiler la variante corriges d'un meme fichier source ; par defaut (non defini),
% la compilation produit la variante enonces. \CAPETSHOW{enonce}{corrige} choisit la branche.
%
% TROIS MODES (refonte 2026-09-20, manuel de revision) :
%   - aucun drapeau   -> imprime l'enonce            (variante enonces, INCHANGE)
%   - \CAPETCORRIGE   -> imprime le corrige          (variante corriges, INCHANGE)
%   - \CAPETMANUEL    -> imprime l'enonce ET met le corrige de cote pour l'annexe
%
% Le mode manuel stocke chaque corrige dans une macro (\csgdef), sans ecriture
% fichier ni re-tokenisation : le contenu (circuitikz, accents, %) traverse intact.
% L'annexe les rejoue ensuite dans l'ordre via \CAPETANNEXE.
\usepackage{etoolbox}

\newcounter{CAPETexo}          % numero continu d'exercice dans le manuel
\setcounter{CAPETexo}{0}
\newcounter{CAPETannexe}       % compteur de relecture en annexe

% Titre courant de la notion en cours (pose par \CAPETNOTION), sert d'intertitre annexe
\newcommand{\CAPETnotioncode}{}
\newcommand{\CAPETnotiontitre}{}

% \CAPETNOTION{code}{titre} : ouvre une notion dans le manuel (memorise le libelle,
% et marque le rang du premier exercice de la notion pour l'intertitre d'annexe).
\newcommand{\CAPETNOTION}[2]{%
  \renewcommand{\CAPETnotioncode}{#1}%
  \renewcommand{\CAPETnotiontitre}{#2}%
  \csgdef{CAPETnotionstart\the\numexpr\value{CAPETexo}+1\relax}{#1 -- #2}%
}

% ── \degree robuste ───────────────────────────────────────────────────
% Certains fichiers d'exercices utilisent \degree. Ni le preambule ni siunitx
% ne le definissent, et \providecommand local ne suffit pas : sous hyperref,
% \degree existe deja comme commande factice (\ERROR), donc \providecommand
% s'abstient et l'appel echoue a la compilation du manuel (erreur invisible
% en compilation isolee). On le (re)definit donc inconditionnellement ici,
% apres le chargement des packages, pour tout le corpus.
\AtBeginDocument{\providecommand{\degree}{}\renewcommand{\degree}{\ensuremath{^\circ}}}

% Libelle decoratif : les fichiers d'exercices utilisent aussi \CAPETSHOW dans
% leur \fancyhead et leur bloc de titre, pour afficher "(enonces)" ou "(corriges)".
% Ces appels-la ne sont PAS des exercices : en mode manuel ils ne doivent ni
% consommer un numero, ni alimenter l'annexe. On les reconnait a leur premier
% argument, strictement egal au libelle ci-dessous.
\newcommand{\CAPETLIBELLEENONCES}{(\'enonc\'es)}

\ifdefined\CAPETMANUEL
  % --- Mode manuel : enonce en place, corrige differe vers l'annexe ---
  \newcommand{\CAPETSHOW}[2]{%
    \def\CAPETargun{#1}%
    \ifx\CAPETargun\CAPETLIBELLEENONCES
      #1%            appel decoratif : rendu tel quel, hors numerotation
    \else
      \CAPETSHOWEXO{#1}{#2}%
    \fi
  }
  \newcommand{\CAPETSHOWEXO}[2]{%
    \stepcounter{CAPETexo}%
    \csgdef{CAPETcorrige\theCAPETexo}{#2}%
    \edef\CAPETtmpnum{\theCAPETexo}%
    #1%
    \par\vspace{2pt}%
    {\footnotesize\itshape\color{colExo}%
     \hyperref[capet:cor:\CAPETtmpnum]{Corrige de l'exercice \CAPETtmpnum\ $\rightarrow$ p.~\pageref{capet:cor:\CAPETtmpnum}}%
     \label{capet:eno:\CAPETtmpnum}}%
  }
  % \CAPETANNEXE : rejoue tous les corriges captures, dans l'ordre, avec intertitres
  \newcommand{\CAPETANNEXE}{%
    \setcounter{CAPETannexe}{0}%
    \loop\ifnum\value{CAPETannexe}<\value{CAPETexo}%
      \stepcounter{CAPETannexe}%
      \edef\CAPETcurnum{\theCAPETannexe}%
      \ifcsname CAPETnotionstart\CAPETcurnum\endcsname
        \par\medskip
        {\normalsize\bfseries\color{colExo}\csuse{CAPETnotionstart\CAPETcurnum}}%
        \par\nobreak\smallskip
      \fi
      \begin{tcolorbox}[exoBox, breakable,
        title={Corrige de l'exercice \CAPETcurnum\
               \hfill\normalfont\footnotesize
               \hyperref[capet:eno:\CAPETcurnum]{enonce p.~\pageref{capet:eno:\CAPETcurnum}}}]%
      \label{capet:cor:\CAPETcurnum}%
      \csuse{CAPETcorrige\CAPETcurnum}%
      \end{tcolorbox}%
    \repeat
  }
\else
  % --- Modes historiques, comportement strictement inchange ---
  \newcommand{\CAPETSHOW}[2]{\ifdefined\CAPETCORRIGE #2\else #1\fi}
  \newcommand{\CAPETANNEXE}{}
\fi

% ── Difficulté 1 à 3 (étoiles pleines, amssymb) ───────────────────────
\newcommand{\difficulte}[1]{%
  \ifcase#1\relax
  \or $\bigstar$\;niveau 1/3%
  \or $\bigstar\bigstar$\;niveau 2/3%
  \or $\bigstar\bigstar\bigstar$\;niveau 3/3%
  \fi
}

% ── Style encadré extrait de document technique (banque d'exercices) ──
\tcbset{
  docTech/.style={
    breakable,
    enhanced,
    colback=colGrayLight,
    colframe=colSeparator,
    fonttitle=\bfseries\small\itshape,
    coltitle=black,
    colbacktitle=colGrayLight,
    left=4pt, right=4pt, top=3pt, bottom=3pt,
    before skip=4pt, after skip=4pt,
  }
}

% ── Style tcolorbox de section principale ────────────────────────────
% Refonte 2026 : bandeau titre coloré arrondi + fond très pâle, sans cadre.
\tcbset{
  ficheSection/.style={
    breakable, enhanced jigsaw,
    frame hidden, boxrule=0pt, arc=1.8mm,
    colback=#1!4!white,
    colbacktitle=#1,
    coltitle=white,
    fonttitle=\sffamily\bfseries,
    left=7pt, right=7pt, top=5pt, bottom=5pt,
    toptitle=3pt, bottomtitle=3pt,
    before skip=10pt, after skip=8pt,
  }
}

% ── Style encadré question type ──────────────────────────────────────
% Carte blanche, bandeau sombre, liseré très clair.
\tcbset{
  questiontype/.style={
    breakable, enhanced jigsaw,
    colback=white,
    frame hidden, boxrule=0pt, arc=1.8mm,
    colbacktitle=colElec!78!black,
    coltitle=white,
    fonttitle=\sffamily\bfseries\small,
    left=7pt, right=7pt, top=4pt, bottom=4pt,
    toptitle=2.5pt, bottomtitle=2.5pt,
    before skip=8pt, after skip=8pt,
    borderline={0.6pt}{0pt}{colElec!12!white},
  }
}

% ── Style encadré jury ───────────────────────────────────────────────
% Bandeau ambre, fond crème, sans cadre.
\tcbset{
  jury/.style={
    breakable, enhanced jigsaw,
    colback=colJuryBg,
    frame hidden, boxrule=0pt, arc=1.8mm,
    colbacktitle=colJuryFrame,
    coltitle=white,
    fonttitle=\sffamily\bfseries,
    left=7pt, right=7pt, top=4pt, bottom=4pt,
    toptitle=3pt, bottomtitle=3pt,
    before skip=10pt, after skip=8pt,
  }
}

% ── En-tête / pied de page communs (refonte 2026) ─────────────────────
% #1 code fiche, #2 titre court, #3 couleur domaine
\newcommand{\ficheHeaders}[3]{%
  \fancyhead[L]{\small\sffamily{\bfseries\textcolor{#3}{#1}}\hspace{0.7em}\textcolor{colSlate}{#2}}%
  \fancyhead[R]{\small\sffamily\textcolor{colSlate}{CAPET SII IE\, --- \,Session 2026}}%
  \fancyfoot[C]{\small\sffamily\textcolor{colSlate}{\thepage\,/\,\pageref{LastPage}}}%
  \fancyhead[C]{}\fancyfoot[L]{}\fancyfoot[R]{}%
}

% ── Bloc en-tête de fiche (BLOC 0, refonte 2026) ──────────────────────
% #1 code  #2 titre  #3 couleur domaine  #4 domaine  #5 ligne présence  #6 sous-thèmes
\newcommand{\ficheEntete}[6]{%
  \begin{tcolorbox}[enhanced, breakable,
    frame hidden, boxrule=0pt, arc=2.5mm,
    interior style={left color=#3, right color=#3!75!black},
    left=9pt, right=9pt, top=7pt, bottom=7pt,
    before skip=0pt, after skip=6pt]
    {\color{white}\sffamily\bfseries\LARGE #1\quad #2}\\[4pt]
    {\color{#3!22!white}\sffamily\small \MakeUppercase{#4}\quad\textbar\quad #5}\\[1pt]
    {\color{#3!35!white}\sffamily\footnotesize Sous-thèmes : #6}
  \end{tcolorbox}%
}

% ── Mise en page ─────────────────────────────────────────────────────
\pagestyle{fancy}
\fancyhf{}
\renewcommand{\headrulewidth}{0.5pt}
\renewcommand{\footrulewidth}{0pt}
\AtBeginDocument{%
  \renewcommand{\headrule}{{\color{colRule}\hrule width\headwidth height 0.5pt}}%
}

\setlist[itemize]{leftmargin=1.2em, itemsep=1pt, topsep=2pt, parsep=0pt}
\setlist[enumerate]{leftmargin=1.2em, itemsep=3pt, topsep=2pt, parsep=0pt}

\setlength{\parskip}{2pt}
\setlength{\columnsep}{10pt}

\def\CAPETBEGINFILE{\begin{document}}
\def\CAPETENDFILE{\end{document}}
\else
\def\CAPETBEGINFILE{}
\def\CAPETENDFILE{}
\fi

\ficheHeaders{T19}{Amplificateur opérationnel}{colElec}

\CAPETBEGINFILE

% -----------------------------------------------------------------------
%  BLOC EN-TETE
% -----------------------------------------------------------------------

\ficheEntete{T19}{AOP (amplificateur opérationnel)}{colElec}%
  {électronique - Conditionnement du signal}%
  {\textbf{Présence :} \textbf{4/13} sujets (2013-2026) - thème \emph{régulier depuis 2019}}%
  {montages inverseur - non-inverseur - suiveur - soustracteur - intégrateur - comparateur - hystérésis - conversion tension-courant 4-20\,mA - théorème de Millman}

% -----------------------------------------------------------------------
%  1. RAPPELS DE COURS
% -----------------------------------------------------------------------

\begin{tcolorbox}[
  ficheSection=colElec,
  title={1. Rappels de cours},
]

% -- 1.1 Hypotheses AOP ideal -----------------------------------------
\textbf{1.1 Hypothèses de l'AOP idéal - régime linéaire}

\smallskip
\begin{center}
\renewcommand{\arraystretch}{1.3}
\begin{tabular}{ll}
\toprule
\textbf{Hypothèse} & \textbf{Conséquence} \\
\midrule
Courants d'entrée nuls & $I_+ = I_- = 0$ \\
Gain différentiel infini & $\varepsilon = V_+ - V_- = 0 \;\Rightarrow\; V_+ = V_-$ \\
\bottomrule
\end{tabular}
\end{center}
\textbf{Condition de validité :} contre-réaction négative (sortie rebouclée sur $V_-$). En comparateur (boucle ouverte ou contre-réaction positive), $V_+ \neq V_-$.

\medskip\hrule\vspace{4pt}

% -- 1.2 Montages fondamentaux ----------------------------------------
\textbf{1.2 Montages fondamentaux}

\smallskip
\begin{center}
\renewcommand{\arraystretch}{1.45}
\begin{tabular}{llll}
\toprule
\textbf{Montage} & \textbf{Gain $V_s/V_e$} & \textbf{$Z_e$} & \textbf{Usage} \\
\midrule
Inverseur & $-R_2/R_1$ & $R_1$ & Amplification + inversion \\
Non-inverseur & $1 + R_2/R_1$ & $\to \infty$ & Amplification sans inversion \\
Suiveur (buffer) & $1$ & $\to \infty$ & Adaptation d'impédance \\
Soustracteur (1 AOP) & $\dfrac{R_2}{R_1}(V_2 - V_1)$ & - & Signal différentiel (capteur) \\
Intégrateur & $-\dfrac{1}{RCp}$ & $R_1$ & Filtre passe-bas, correcteur I \\
\bottomrule
\end{tabular}
\end{center}

\smallskip
\textbf{Attention - deux topologies différentielles distinctes.} Le soustracteur ci-dessus (1 seul AOP, quatre résistances appariées) donne $V_s = (R_2/R_1)(V_2 - V_1)$. L'\emph{étage d'entrée d'un amplificateur d'instrumentation} (2 AOP rebouclés par $R_2$ entre les deux entrées inverseuses) a un gain différent : $G_1 = \dfrac{2R_1 + R_2}{R_2}$ (structure du sujet 2024, cf.\ QT1). Ne pas mélanger les deux formules.

\medskip\hrule\vspace{4pt}

% -- 1.3 Comparateur - hysteresis -------------------------------------
\textbf{1.3 Comparateur à hystérésis - contre-réaction positive}

\smallskip
Contre-réaction sur $V_+$ (résistances $R_1$ sur sortie, $R_2$ sur $V_{ref}$) :
\[
  V_+ = \frac{R_2\,V_{ref} + R_1\,V_s}{R_1 + R_2}
\]
Deux seuils de basculement (avec $V_s = \pm V_{sat}$) :
\[
  V_H = \frac{R_2\,V_{ref} + R_1\,V_{sat}}{R_1 + R_2}
  \qquad;\qquad
  V_B = \frac{R_2\,V_{ref} - R_1\,V_{sat}}{R_1 + R_2}
\]

\smallskip
\textbf{Cycle d'hystérésis :}
\begin{center}
\begin{tikzpicture}[scale=0.9, font=\footnotesize]
  \draw[->] (0,0) -- (5.5,0) node[right] {$V_+$};
  \draw[->] (0,-1.8) -- (0,1.8) node[above] {$V_s$};
  \draw[thick, colElec] (0.3,-1.3) -- (3.0,-1.3);
  \draw[thick, colElec, ->] (3.0,-1.3) -- (3.0,1.3);
  \draw[thick, colElec] (3.0,1.3) -- (4.8,1.3);
  \draw[thick, colElec!55] (1.3,1.3) -- (3.0,1.3);
  \draw[thick, colElec!55, ->] (1.3,1.3) -- (1.3,-1.3);
  \draw[thick, colElec!55] (0.3,-1.3) -- (1.3,-1.3);
  \draw[dashed] (1.3,0) -- (1.3,-1.3);
  \node[below] at (1.3,0) {\scriptsize $V_B$};
  \draw[dashed] (3.0,0) -- (3.0,1.3);
  \node[below] at (3.0,0) {\scriptsize $V_H$};
  \node[left] at (0,1.3) {\scriptsize $+V_{sat}$};
  \node[left] at (0,-1.3) {\scriptsize $-V_{sat}$};
\end{tikzpicture}
\end{center}
Trajet foncé : $V_+$ croissant, basculement à $V_H$. Trajet clair : $V_+$ décroissant, basculement à $V_B$. L'écart $V_H - V_B$ évite les oscillations parasites au voisinage du seuil.

\medskip\hrule\vspace{4pt}

% -- 1.4 Conversion tension-courant -----------------------------------
\textbf{1.4 Conversion tension→courant 4-20\,mA (transmetteur normalisé)}

\smallskip
\[
  I_s = \frac{V_{entree}}{R_5}
\]
Pour $V_{entree} \in [1\,;\,5]\,\mathrm{V}$ et $R_5 = 250\,\Omega$ : $I_s \in [4\,;\,20]\,\mathrm{mA}$.\\
$I_s$ est indépendant de la charge $R_{ana}$ car la boucle de contre-réaction maintient $\varepsilon = 0$.

\textbf{Théorème de Millman} (pour un nœud à $n$ admittances) :
\[
  V_{noeud} = \frac{\sum_k V_k/R_k}{\sum_k 1/R_k}
\]

\end{tcolorbox}

% -----------------------------------------------------------------------
%  2. FORMULAIRE SYNTHETIQUE
% -----------------------------------------------------------------------

\begin{tcolorbox}[
  ficheSection=colElec,
  title={2. Formulaire synthétique},
]
\small
\begin{multicols}{2}
\renewcommand{\arraystretch}{1.45}
\begin{tabular}{@{}p{4.5cm}ll@{}}
\toprule
\textbf{Formule} & \textbf{Usage} & \textbf{U.} \\
\midrule
\multicolumn{3}{@{}l}{\footnotesize\textit{Montages linéaires}} \\
$V_s/V_e = -R_2/R_1$ & Inverseur & - \\
$V_s/V_e = 1+R_2/R_1$ & Non-inv. & - \\
$V_s = V_e$ & Suiveur & - \\
$V_s = (R_2/R_1)(V_2-V_1)$ & Soustr. 1 AOP & - \\
$G_1 = (2R_1+R_2)/R_2$ & Instrum. 2 AOP & - \\
$H(p) = -1/(RCp)$ & Intégr. & - \\
\midrule
\multicolumn{3}{@{}l}{\footnotesize\textit{Hypothèses}} \\
$V_+ = V_-$ & Lin. & - \\
$I_+ = I_- = 0$ & Lin. & - \\
\bottomrule
\end{tabular}

\columnbreak

\renewcommand{\arraystretch}{1.45}
\begin{tabular}{@{}p{4.5cm}ll@{}}
\toprule
\textbf{Formule} & \textbf{Usage} & \textbf{U.} \\
\midrule
\multicolumn{3}{@{}l}{\footnotesize\textit{Hystérésis}} \\
$V_+ = (R_2 V_{ref}+R_1 V_s)/(R_1+R_2)$ & Seuil & V \\
$V_H = (R_2 V_{ref}+R_1 V_{sat})/(R_1+R_2)$ & Haut & V \\
$V_B = (R_2 V_{ref}-R_1 V_{sat})/(R_1+R_2)$ & Bas & V \\
\midrule
\multicolumn{3}{@{}l}{\footnotesize\textit{T/I et Millman}} \\
$I_s = V_e/R_5$ & Transm. & A \\
$V_n = \sum(V_k/R_k)/\sum(1/R_k)$ & Millman & V \\
\midrule
\multicolumn{3}{@{}l}{\footnotesize\textit{Données 2024}} \\
$G_{total} = 340$ & Transmetteur & - \\
Plage : 4-20\,mA, $R_5 = 250\,\Omega$ & Normalisé & - \\
\bottomrule
\end{tabular}
\end{multicols}
\normalsize
\end{tcolorbox}

% -----------------------------------------------------------------------
%  3. QUESTIONS TYPES
% -----------------------------------------------------------------------

\begin{tcolorbox}[
  ficheSection=colElec,
  title={3. Questions types - extraites des sujets réels},
]

% -- QT1 -------------------------------------------------------------

\begin{tcolorbox}[questiontype,
  title={QT1 - Étage différentiel d'instrumentation - calcul du gain $G_1$ (transmetteur de température) \hfill\normalfont\footnotesize E2 disc. 2024 (E1-E3)}
]
\textbf{Énoncé.}
Un étage d'entrée différentiel à 2 AOP (structure d'amplificateur d'instrumentation, distincte du soustracteur à 1 AOP du cours 1.2) amplifie la tension $e_{AB}$ d'un capteur PT100. Données : $R_1 = 1\,\mathrm{k\Omega}$, $R_2 = 834{,}4\,\Omega$. Exprimer le gain $G_1$ et calculer sa valeur. Le 2\textsuperscript{e} étage a $G_2 = -100$, le 3\textsuperscript{e} a $G_3 = -1$. Calculer le gain global.

\tcbline

\textbf{Corrigé.}

\textbf{Étape 1 — Gain du 1er étage.} Par la loi des mailles (AOP idéaux, $V_+ = V_-$, courant commun dans $R_1$--$R_2$--$R_1$), le gain de l'étage d'entrée à 2 AOP vaut :
\[
  G_1 = \frac{2R_1 + R_2}{R_2} = \frac{2 \times 1000 + 834{,}4}{834{,}4} = \frac{2834{,}4}{834{,}4} \approx \mathbf{3{,}4}
\]

\textbf{Étape 2 — Gain global.}
\[
  G = G_1 \times G_2 \times G_3 = 3{,}4 \times (-100) \times (-1) = \mathbf{340}
\]

\textbf{Conclusion.} $V_3 = 340 \times e_{AB}$ : les deux inversions de $G_2$ et $G_3$ se compensent, le signal final reste en phase avec $e_{AB}$, ce qui simplifie l'exploitation en aval (pas de correction de signe nécessaire).
\end{tcolorbox}

\vspace{4pt}

% -- QT2 -------------------------------------------------------------

\begin{tcolorbox}[questiontype,
  title={QT2 - Comparateur à hystérésis - seuils de basculement \hfill\normalfont\footnotesize E2 disc. 2023 (F6-F8)}
]
\textbf{Énoncé.}
Comparateur à hystérésis : $V_{ref} = 2{,}42\,\mathrm{V}$, $V_{sat} = 15\,\mathrm{V}$, $R_1 = 1\,\mathrm{k\Omega}$, $R_2 = 150\,\mathrm{k\Omega}$.
(a) Exprimer $V_+$ en fonction de $V_{ref}$, $R_1$, $R_2$, $V_s$.
(b) Calculer $V_H$ et $V_B$.

\tcbline

\textbf{Corrigé.}

\textbf{Étape 1 (a) — Diviseur de tension entre $V_s$ et $V_{ref}$.}
\[
  V_+ = \frac{R_2\,V_{ref} + R_1\,V_s}{R_1 + R_2} = \frac{150 \times 10^3 \times V_{ref} + 1 \times 10^3 \times V_s}{151 \times 10^3}
\]

\textbf{Étape 2 (b) — Seuil haut} ($V_s = +V_{sat} = +15\,\mathrm{V}$).
\[
  V_H = \frac{150\,000 \times 2{,}42 + 1\,000 \times 15}{151\,000} = \frac{363\,000 + 15\,000}{151\,000} = \frac{378\,000}{151\,000} \approx \mathbf{2{,}50\,\mathrm{V}}
\]

Seuil bas ($V_s = -V_{sat} = -15\,\mathrm{V}$) :
\[
  V_B = \frac{150\,000 \times 2{,}42 - 1\,000 \times 15}{151\,000} = \frac{363\,000 - 15\,000}{151\,000} = \frac{348\,000}{151\,000} \approx \mathbf{2{,}30\,\mathrm{V}}
\]

\textbf{Conclusion.} L'hystérésis $\Delta V = V_H - V_B = 0{,}20\,\mathrm{V}$ crée une zone morte autour du seuil nominal : c'est cette marge, et non un seuil unique, qui protège le montage contre les basculements intempestifs dus au bruit de mesure.
\end{tcolorbox}

\vspace{4pt}

% -- QT3 -------------------------------------------------------------

\begin{tcolorbox}[questiontype,
  title={QT3 - Conversion tension→courant - indépendance vis-à-vis de la charge \hfill\normalfont\footnotesize E2 disc. 2024 (E9)}
]
\textbf{Énoncé.}
Montage AOP en boucle fermée avec résistance de mesure $R_5 = 250\,\Omega$ dans la boucle. Tension en entrée $V_4 \in [1\,;\,5]\,\mathrm{V}$. Montrer que le courant $I_s$ est indépendant de la charge $R_{ana}$ et calculer la plage.

\tcbline

\textbf{Corrigé.}

\textbf{Étape 1 — Loi des nœuds au nœud A} (entrée $V_-$ de l'AOP, hypothèse $I_- = 0$) :
\[
  \frac{V_A - V_4}{R_5} + \frac{V_A - V_s}{R_{ana}} = I_- = 0
\]
Avec $V_+ = V_- = V_A$ (contre-réaction), et en appliquant la règle d'entrée $V_- = V_4$ :
\[
  I_s = \frac{V_4 - 0}{R_5} = \frac{V_4}{R_5}
\]
$I_s$ est bien \textbf{indépendant de $R_{ana}$} (la boucle de contre-réaction ajuste $V_s$ pour maintenir $\varepsilon = 0$ quelle que soit la charge).

\textbf{Étape 2 — Plage de courant.}
\[
  V_4 = 1\,\mathrm{V} \;\Rightarrow\; I_s = \frac{1}{250} = 4\,\mathrm{mA}
  \qquad;\qquad
  V_4 = 5\,\mathrm{V} \;\Rightarrow\; I_s = \frac{5}{250} = 20\,\mathrm{mA}
\]

\textbf{Conclusion.} La plage obtenue (4-20\,mA) correspond exactement au standard industriel normalisé : ce montage est directement exploitable comme sortie de transmetteur, sans étage de conversion supplémentaire.
\end{tcolorbox}

\end{tcolorbox}

% -----------------------------------------------------------------------
%  4. ATTENTES DU JURY
% -----------------------------------------------------------------------

\begin{tcolorbox}[
  jury,
  title={4. Attentes du jury - Erreurs récurrentes et points de vigilance},
]

\begin{enumerate}[label=\textbf{\arabic*.}, leftmargin=*, itemsep=5pt]

  \item \textbf{Vérifier la contre-réaction avant d'appliquer $V_+ = V_-$.}
  En mode comparateur (boucle ouverte ou contre-réaction positive), $V_+ \neq V_-$ - ne pas appliquer l'égalité. Identifier d'abord le mode de fonctionnement.

  \item \textbf{Signe du gain inverseur.}
  $G = -R_2/R_1$ : le signe moins est fondamental. L'oubli conduit à conclure à une amplification non-inverseuse.

  \item \textbf{Courants d'entrée AOP nuls.}
  $I_+ = I_- = 0$ : appliquer les lois de nœuds uniquement sur les résistances externes. Ne jamais écrire un courant passant par l'entrée de l'AOP.

  \item \textbf{Théorème de Millman méconnu.}
  Appliqué en 2024 (E4 : montage de décalage), il simplifie considérablement le calcul des nœuds à plusieurs branches.

  \item \textbf{Hystérésis : ne pas inverser les seuils.}
  $V_H > V_B$ toujours. $V_H$ correspond à $V_s = +V_{sat}$, $V_B$ à $V_s = -V_{sat}$. Confondre les deux donne une hystérésis négative physiquement impossible.

  \item \textbf{Suiveur = adaptation d'impédance, pas amplification.}
  Gain = 1 mais impédance d'entrée quasi infinie et impédance de sortie quasi nulle. Rôle : isoler deux étages sans atténuation.

\end{enumerate}

\medskip
\begin{tcolorbox}[
  colback=colJuryFrame!12!colJury,
  colframe=colJuryFrame,
  left=4pt, right=4pt, top=2pt, bottom=2pt,
  fontupper=\small\itshape
]
Rapport jury 2024, sur la partie du sujet consacrée à l'électronique analogique : «~la cinquième et dernière partie a été abordée par 70\,\% des candidats et 47\,\% d'entre eux l'ont traitée correctement.~»\\
Rapport jury 2024 : «~Le jury attend des réponses détaillées et argumentées. Des résultats donnés directement, sans calcul et sans justification ne répondent pas aux attendus de l'épreuve.~»
\end{tcolorbox}

\end{tcolorbox}

% -----------------------------------------------------------------------
%  PIED - PRESENCE DANS LES SUJETS
% -----------------------------------------------------------------------

\begin{tcolorbox}[
  ficheSection=colPresence,
  title={Présence dans les sujets analysés (2013-2026)},
]
\begin{center}
\footnotesize\renewcommand{\arraystretch}{1.15}
\begin{tabular}{lccccccccccccc}
\toprule
\textbf{Année} & 13 & 14 & 15 & 16 & 17 & 18 & 19 & 20 & 21 & 22 & 23 & 24 & 25 \\
\midrule
E1 (commune)   & - & - & - & - & - & - & - & - & - & - & - & $\checkmark$ & - \\
E2 discipl. IE & - & - & - & - & - & - & $\checkmark$ & - & - & $\checkmark$ & $\checkmark$ & $\checkmark$ & - \\
E2 appl. comm. & - & - & - & - & - & - & - & - & - & - & - & - & - \\
\bottomrule
\end{tabular}
\end{center}
\smallskip
\footnotesize\textbf{Sous-thèmes fréquents :}
montage inverseur/non-inverseur (6/13) - montage soustracteur (4/13) - comparateur à hystérésis (3/13) - conversion tension-courant 4-20\,mA (2/13) - intégrateur/correcteur PI (3/13)

\medskip
\textbf{Ordres de grandeur (repères rapides).}
\begin{itemize}[leftmargin=1.3em]
  \item Réseau BT FR : tension phase-neutre typique \SI{230}{\volt}, tension phase-phase typique \SI{400}{\volt}, fréquence \SI{50}{\hertz}.
  \item Courants d'étude fréquents en sujet : de quelques \si{\ampere} à quelques dizaines de \si{\ampere}.
  \item En conversion d'énergie, un rendement attendu reste généralement entre 0{,}8 et 0{,}98 selon la technologie.
\end{itemize}

\medskip
\textbf{Synthèse express (à refaire sans support).}
\begin{enumerate}[leftmargin=1.4em]
  \item Poser les hypothèses du modèle et les conventions de signe.
  \item Écrire 2 à 4 relations clés avec unités.
  \item Réaliser un calcul numérique complet et vérifier l'ordre de grandeur.
  \item Justifier le résultat par un argument technique de niveau CAPET.
\end{enumerate}

\end{tcolorbox}

\CAPETENDFILE
